\documentclass[11pt]{article}

\usepackage[margin=0.9in]{geometry}
\usepackage{amsmath,amssymb}
\usepackage{booktabs}
\usepackage{enumitem}
\usepackage{listings}
\usepackage{xcolor}
\usepackage{hyperref}
\usepackage{tikz}

\definecolor{backcolour}{rgb}{0.95,0.95,0.92}
\definecolor{codeblue}{rgb}{0.05,0.15,0.55}
\definecolor{codegreen}{rgb}{0.0,0.45,0.0}

\lstset{
    backgroundcolor=\color{backcolour},
    basicstyle=\ttfamily\small,
    keywordstyle=\color{codeblue}\bfseries,
    commentstyle=\color{codegreen},
    stringstyle=\color{purple},
    breaklines=true,
    frame=single,
    showstringspaces=false,
    tabsize=4
}

\title{CPSC 250 \\ Derived Classes and Default Values}
\author{Edward Brash}
\date{}

\begin{document}

\maketitle

\section{Motivation}

Object-oriented programming becomes much more powerful when we can create new classes from existing classes.

This is called \textbf{inheritance}.

A derived class inherits variables and methods from a base class, then adds more specific behavior.

\section{Base Class: Transportation}

Suppose we begin with a general class representing a transportation mode.

\begin{lstlisting}[language=Python]
class Transportation:

    def __init__(self, name, speed=75):
        self.name = name
        self.speed = speed

    def info(self):
        print(f"{self.name} can go {self.speed} mph")
\end{lstlisting}

The default speed is:

\[
75 \text{ mph}
\]

This means that if the user does not specify a speed, Python uses 75 automatically.

\section{Why Use a Base Class?}

There are many modes of transportation:

\begin{itemize}
    \item motorized vehicles,
    \item self-powered vehicles,
    \item flight vehicles,
    \item water vehicles,
    \item and many others.
\end{itemize}

All of these are transportation modes, so it is useful to put common features in a shared base class.

\section{Derived Class: MotorizedVehicle}

A motorized vehicle is a transportation mode, but it has additional properties.

For example, it has:

\begin{itemize}
    \item miles per gallon,
    \item fuel level,
    \item ability to add fuel,
    \item ability to drive.
\end{itemize}

\begin{lstlisting}[language=Python]
class MotorizedVehicle(Transportation):

    def __init__(self, name, speed=80, mpg=40):
        Transportation.__init__(self, name, speed)

        self.mpg = mpg
        self.fuel_gal = 0

    def add_fuel(self, amount):
        self.fuel_gal += amount

    def drive(self, distance):
        required_fuel = distance / self.mpg

        if self.fuel_gal < required_fuel:
            print("Not enough fuel")

        else:
            self.fuel_gal -= required_fuel
            print(f"{self.fuel_gal} gallons remain")
\end{lstlisting}

\section{Explicitly Calling the Base Constructor}

The line:

\begin{lstlisting}[language=Python]
Transportation.__init__(self, name, speed)
\end{lstlisting}

explicitly calls the constructor of the base class.

This initializes the part of the object inherited from \texttt{Transportation}.

This is extremely important: the derived class constructor does not automatically run the base class constructor unless we ask it to.

\section{Derived Class: Motorcycle}

A motorcycle is a specific kind of motorized vehicle.

\begin{lstlisting}[language=Python]
class Motorcycle(MotorizedVehicle):

    def __init__(self, name, speed=55, mpg=25):
        MotorizedVehicle.__init__(self, name, speed, mpg)

    def wheelie(self):
        print("That's dangerous!")
\end{lstlisting}

The \texttt{Motorcycle} class inherits:

\begin{itemize}
    \item \texttt{name},
    \item \texttt{speed},
    \item \texttt{mpg},
    \item \texttt{fuel\_gal},
    \item \texttt{info()},
    \item \texttt{add\_fuel()},
    \item \texttt{drive()}.
\end{itemize}

It also adds:

\begin{lstlisting}[language=Python]
wheelie()
\end{lstlisting}

\section{Using the Motorcycle Class}

\begin{lstlisting}[language=Python]
scooter = Motorcycle("Vespa", 55, 40)
dirtbike = Motorcycle("KX450F", 80, 25)

scooter.info()
dirtbike.info()
\end{lstlisting}

We can also rely on default values:

\begin{lstlisting}[language=Python]
harley = Motorcycle("Harley Davidson")

harley.info()
print(harley.mpg)
print(harley.fuel_gal)
\end{lstlisting}

The object \texttt{harley} is created using the default values from the \texttt{Motorcycle} constructor.

\section{A Generic Motorized Vehicle}

We can still instantiate the parent class directly:

\begin{lstlisting}[language=Python]
generic = MotorizedVehicle("People Carrier")

generic.info()
print(generic.mpg)
print(generic.fuel_gal)
\end{lstlisting}

This is useful when the base or intermediate class is meaningful on its own.

\section{Inheritance Chain}

The inheritance chain in this example is:

\[
\texttt{Transportation}
\rightarrow
\texttt{MotorizedVehicle}
\rightarrow
\texttt{Motorcycle}
\]

Each derived class inherits everything from the classes above it, unless behavior is overridden or hidden.

\section{Key Takeaways}

\begin{itemize}
    \item A derived class inherits from a base class.
    \item Default values make constructors more flexible.
    \item Derived constructors should initialize the base part of the object.
    \item In Python, this is often done explicitly with a base-class constructor call.
    \item Inheritance allows us to reuse common code while adding specialized behavior.
\end{itemize}

\end{document}
